% NNGraph DSL -- English LaTeX Report
% Compile: xelatex -shell-escape report-en.tex  (run twice)

\documentclass{nngraph-en}

\usepackage{fontspec}
\setmonofont{Liberation Mono}

% ----------------------------------------------------------------
\title{
    \textbf{NNGraph DSL Compiler} \\[0.5em]
    \large A Domain-Specific Language for Describing Neural Network Architectures \\
    and Automatic PyTorch Code Generation
}
\author{\large Project Technical Report}
\date{\today}

% ----------------------------------------------------------------
\begin{document}

\maketitle
\tableofcontents
\newpage

% ================================================================
\chapter*{Abstract}
\markboth{Abstract}{Abstract}
\addcontentsline{toc}{chapter}{Abstract}

NNGraph DSL is a domain-specific language designed to describe
neural network architectures as directed graphs. It accepts
\code{.nng} files as input and produces executable PyTorch code
consisting of a complete \code{nn.Module} subclass with
\code{\_\_init\_\_} and \code{forward} methods.

The compiler is built with ANTLR4 version 4.13.2 and uses the
Listener pattern for semantic analysis. Kahn's topological sort
algorithm is employed both in the semantic analyzer (cycle detection)
and in the code generator (execution ordering). Tensor shape
inference, Graphviz graph visualization, training scaffolding,
quantization annotations, and ONNX export are also supported.

Four example
architectures---MLP, ResBlock, InceptionBlock, and
TransformerEncoder---compile successfully and produce correct
output tensor shapes in their forward pass.

% ================================================================
\chapter{Introduction}

\section{Motivation}

Writing \code{nn.Module} subclasses in PyTorch by hand for complex
architectures is repetitive and error-prone. For a network with
ten layers, the programmer must:

\begin{itemize}
    \item Define each layer separately with correct parameters in
    \code{\_\_init\_\_}
    \item Manually manage the execution order in \code{forward}
    \item Name intermediate variables in branching topologies
    \item Ensure no shared tensor is overwritten before it is consumed
\end{itemize}

Consider the following PyTorch code that must be written manually:

\begin{codebox}{Manual boilerplate for a ResBlock}{code-resblock-manual}
\begin{amzcode}{python}
class ResBlock(nn.Module):
    def __init__(self):
        super(ResBlock, self).__init__()
        self.conv1 = nn.Conv2d(
            in_channels=64, out_channels=64,
            kernel_size=3, stride=1, padding=1)
        self.bn1   = nn.BatchNorm2d(num_features=64)
        self.relu1 = nn.ReLU()
        # ... and so on for each layer

    def forward(self, x):
        identity = x
        x = self.conv1(x)
        x = self.bn1(x)
        x = self.relu1(x)
        # ... manage residual connection manually
        x = x + identity
        return x
\end{amzcode}
\end{codebox}

With NNGraph DSL, the architecture is described once as a graph and
the compiler generates this code automatically.

\section{Goal}

The goal of NNGraph DSL is to make the \code{.nng} file the single
source of truth for the architecture. The programmer describes the
graph; the compiler generates all boilerplate.

\section{Scope}

This tool is suited for describing static architectures---networks
whose computational graph is known at model definition time.
Support for dynamic control flow within \code{forward} is planned
for future versions.

% ================================================================
\chapter{Related Work}

\section{The NADER Framework}

\begin{dfnbox}{NADER}{nader}
\dfntxt{NADER} (Neural Architecture Design via Representation) is a
framework that models neural network architectures as directed
computational graphs (DAGs). Each node represents a layer operation
and each edge represents tensor flow.
\end{dfnbox}

NNGraph DSL adopts this graph model as the foundation of its language
design.

\section{ANTLR4}

\begin{dfnbox}{ANTLR4}{antlr4}
\dfntxt{ANTLR4} (ANother Tool for Language Recognition) is a tool
that automatically generates a lexer, parser, and listener/visitor
from a \code{.g4} grammar file.
Version used: ANTLR 4.13.2.
\end{dfnbox}

\subsection{Listener Pattern vs.\ Visitor Pattern}

ANTLR4 provides two patterns for traversing the parse tree:

\begin{itemize}
    \item \textbf{Listener}: \code{ParseTreeWalker} automatically
    traverses the tree and calls \code{enter*/exit*} methods. No
    return values are needed.
    \item \textbf{Visitor}: The programmer controls traversal and
    each method returns a value.
\end{itemize}

\begin{notebox}
This compiler uses the Listener pattern. The \code{exit*} methods
in \code{custom\_listener.py} accumulate state into \code{self.nodes}
and \code{self.edges}---this approach is more natural for building
a graph model from a parse tree.
\end{notebox}

\section{AST Utilities and Visualizations}

This compiler incorporates abstract syntax tree (AST) utilities to support optional structure inspections and visualizations. Three helper files are used:

\begin{itemize}
    \item \code{ast.py} --- \code{AST} and \code{TreeNode} classes
    \item \code{make\_ast\_subtree.py} --- parse tree node composition
    \item \code{ast\_to\_networkx\_graph.py} --- optional graph
    visualization
\end{itemize}

\subsection{Why AST Is Not Used Here}

Most traditional compilers represent programs as abstract syntax trees
(ASTs). Expression-based languages---where operators have precedence
and scopes nest hierarchically---naturally map to tree structures.

NNGraph DSL, however, describes neural network architectures as
\emph{directed acyclic graphs} (DAGs). Nodes represent layer operations;
edges represent tensor flow. The graph structure is the domain model,
not the parse tree.

The compiler performs multiple graph-based operations that require
adjacency lists and explicit edge relationships:

\begin{itemize}[noitemsep]
    \item \textbf{Topological sort} (Kahn's algorithm) for execution
    ordering in \code{forward()}
    \item \textbf{Cycle detection} via topological sort (DAG validation)
    \item \textbf{Reachability analysis} via BFS from input to output
    \item \textbf{Residual connection validation} (checking in-degree
    of merge nodes)
    \item \textbf{Shape inference} by propagating tensor shapes along
    edges
    \item \textbf{Graphviz DOT export} for architecture visualization
    \item \textbf{ONNX export} (ONNX uses a graph format)
\end{itemize}

An AST would encode the parse tree hierarchy (\code{model\_block},
\code{graph\_block}, \code{node\_decl}), but it would obscure the
computational graph that the language is designed to express. For this
reason, the semantic analyzer stores nodes and edges directly in
dictionaries and lists---the natural representation for graph algorithms.

\begin{notebox}
The AST construction code (\code{exitEveryRule}, \code{exitStart})
remains in \code{custom\_listener.py} for optional debugging
visualization via \code{-{}-show-ast}. It is not used during
compilation.
\end{notebox}

% ================================================================
\chapter{Language Design}

\section{Program Structure}

Every \code{.nng} file consists of three blocks. The \code{config}
block is optional:

\begin{codebox}{General structure of a .nng file}{code-structure}
\begin{amzcode}{text}
model <Name> {
    input  <id> : tensor(<shape>)
    output <id>
}

graph {
    node <id> : <LayerType>(<params>)
    ...
    edge <src> -> <dst>
    edge <src> -> <dst> [label="<string>"]
    ...
}

config {
    batch_size = <int>
    device     = "cpu" | "cuda"
    loss       = "<LossFn>"       // enables training scaffold
    optimizer  = "<Optimizer>"    // Adam | SGD | AdamW | RMSprop
    lr         = <float>          // learning rate
    epochs     = <int>            // training epochs
}
\end{amzcode}
\end{codebox}

\begin{dfnbox}{Model block}{block-model}
The \code{model} block specifies the generated class name, the input
node identifier and tensor shape, and the output node identifier.
\end{dfnbox}

\begin{dfnbox}{Graph block}{block-graph}
The \code{graph} block defines all layer nodes and directed edges.
Edge declaration order does not matter; the compiler computes
execution order from the graph structure.
\end{dfnbox}

\begin{dfnbox}{Config block}{block-config}
The \code{config} block sets values for the \code{\_\_main\_\_}
block of the generated code. \code{batch\_size} defaults to 1 and
\code{device} defaults to \code{'cpu'}. If the \code{loss} key is
set, the code generator produces a complete training loop with
optimizer and loss function.
\end{dfnbox}

\section{Grammar}

\begin{dfnbox}{NNGraph.g4 statistics}{grammar-stats}
\begin{itemize}[noitemsep]
    \item Grammar file: \code{NNGraph.g4} (56 lines)
    \item Parser rules: 14 rules (snake\_case)
    \item Lexer tokens: 20 tokens (ALL\_CAPS)
\end{itemize}
\end{dfnbox}

Complete grammar rules:

\begin{codebox}{NNGraph.g4 grammar}{code-grammar}
\begin{amzcode}{antlr}
grammar NNGraph;
start     : model_block graph_block config_block? EOF;
model_block : MODEL ID '{' input_decl output_decl '}';
input_decl  : INPUT ID ':' TENSOR shape_expr;
output_decl : OUTPUT ID;
graph_block : GRAPH '{' (node_decl | edge_decl)* '}';
node_decl   : NODE ID ':' layer_expr;
layer_expr  : ID '(' param_list? ')';
param_list  : param (',' param)*;
param       : ID '=' value;
edge_decl   : EDGE ID ARROW ID ('[' LABEL '=' STRING ']')?;
config_block: CONFIG '{' config_entry* '}';
config_entry: ID '=' value;
value       : FLOAT_LITERAL | INT_LITERAL | BOOL_LITERAL
            | STRING | NONE | shape_expr;
shape_expr  : '(' INT_LITERAL (',' INT_LITERAL)* ')';
\end{amzcode}
\end{codebox}

\begin{notebox}
\textbf{Important design note:} \code{input\_decl} uses
\code{TENSOR shape\_expr} rather than
\code{TENSOR '(' shape\_expr ')'} because \code{shape\_expr}
itself includes the parentheses. This common mistake in early
grammar versions was corrected.
\end{notebox}

\section{Type System}

\begin{center}
\begin{tabular}{lll}
\toprule
\textbf{DSL Type} & \textbf{Python Type} & \textbf{Example} \\
\midrule
\code{int}    & \code{int}        & \code{128} \\
\code{float}  & \code{float}      & \code{0.1} \\
\code{bool}   & \code{bool}       & \code{true} / \code{false} \\
\code{string} & \code{str}        & \code{"relu"} \\
\code{shape}  & \code{tuple[int]} & \code{(64, 32, 32)} \\
\code{None}   & \code{NoneType}   & \code{None} \\
\bottomrule
\end{tabular}
\end{center}

The distinction between \code{int} and \code{float} is enforced at
compile time. For example, \code{Dropout(p=1)} produces an error
because \code{p} expects a \code{float}, not an \code{int}.

\section{Layer Types}

\subsection{Parameterized Layers}

These layers generate a \code{self.<id> = nn.X(...)} line in
\code{\_\_init\_\_}.

\begin{center}
\begin{tabular}{llp{7.5cm}}
\toprule
\textbf{DSL Keyword} & \textbf{PyTorch Class} & \textbf{Key Parameters} \\
\midrule
\code{Linear}        & \code{nn.Linear}             & \code{in\_features}, \code{out\_features}, \code{bias} \\
\code{Conv2d}        & \code{nn.Conv2d}             & \code{in\_ch}, \code{out\_ch}, \code{kernel}, \code{stride}, \code{padding} \\
\code{Conv1d}        & \code{nn.Conv1d}             & \code{in\_ch}, \code{out\_ch}, \code{kernel}, \code{stride} \\
\code{BatchNorm2d}   & \code{nn.BatchNorm2d}        & \code{num\_features} \\
\code{LayerNorm}     & \code{nn.LayerNorm}          & \code{normalized\_shape} \\
\code{MaxPool2d}     & \code{nn.MaxPool2d}          & \code{kernel}, \code{stride} \\
\code{AvgPool2d}     & \code{nn.AvgPool2d}          & \code{kernel}, \code{stride} \\
\code{Dropout}       & \code{nn.Dropout}            & \code{p} \\
\code{Flatten}       & \code{nn.Flatten}            & \code{start\_dim}, \code{end\_dim} \\
\code{Embedding}     & \code{nn.Embedding}          & \code{num\_embeddings}, \code{embedding\_dim} \\
\code{MultiHeadAttn} & \code{nn.MultiheadAttention} & \code{embed\_dim}, \code{num\_heads} \\
\code{LSTM}          & \code{nn.LSTM}               & \code{input\_size}, \code{hidden\_size}, \code{num\_layers} \\
\code{GRU}           & \code{nn.GRU}                & \code{input\_size}, \code{hidden\_size} \\
\bottomrule
\end{tabular}
\end{center}

\subsection{Activations}

\code{ReLU}, \code{Sigmoid}, \code{Tanh}, \code{GELU},
\code{Softmax(dim)}, \code{LeakyReLU(negative\_slope)},
\code{ELU(alpha)}.

\subsection{Graph Operations (No Module)}

These types do not generate a line in \code{\_\_init\_\_}:

\begin{center}
\begin{tabular}{llp{7.5cm}}
\toprule
\textbf{Keyword} & \textbf{Parameters} & \textbf{Generated Code} \\
\midrule
\code{Add}      & ---                          & \code{out = a + b} \\
\code{Concat}   & \code{dim: int}              & \code{out = torch.cat([a,b], dim=d)} \\
\code{Residual} & ---                          & \code{out = main + shortcut} \\
\code{Split}    & \code{chunks: int, dim: int} & \code{split\_0, split\_1, ... = torch.chunk(x,n,dim=d)} \\
\bottomrule
\end{tabular}
\end{center}

\subsection{Universal Parameters}

The \code{quant} parameter can be applied to any node regardless
of layer type. Valid values: \code{"dynamic"}, \code{"static"},
\code{"qat"}.

\begin{codebox}{Quantization annotation example}{code-quant-example}
\begin{amzcode}{text}
node fc1 : Linear(in_features=784, out_features=128, quant="dynamic")
\end{amzcode}
\end{codebox}

\section{Edge Syntax}

\begin{dfnbox}{Simple edge}{edge-simple}
\texttt{edge src -> dst}

Defines a directed edge from node \code{src} to node \code{dst}.
\end{dfnbox}

\begin{dfnbox}{Labeled edge}{edge-labeled}
\texttt{edge src -> dst [label="name"]}

Defines an edge with a string label. Labels serve two roles:
\begin{itemize}
    \item \code{"shortcut"}: marks the skip connection edge into
    a \code{Residual()} node.
    \item \code{"residual\_*"}: marks a skip connection into a
    node with two inputs (Transformer pattern).
\end{itemize}
\end{dfnbox}

\section{Complete Example: MLP}

The \code{.nng} file for a three-layer network:

\begin{codebox}{mlp.nng}{code-mlp-full}
\begin{amzcode}{text}
model MLP {
    input  x   : tensor(784)
    output out
}

graph {
    node fc1   : Linear(in_features=784, out_features=256)
    node relu1 : ReLU()
    node fc2   : Linear(in_features=256, out_features=128)
    node relu2 : ReLU()
    node drop  : Dropout(p=0.3)
    node fc3   : Linear(in_features=128, out_features=10)
    node out   : Softmax(dim=1)

    edge x     -> fc1
    edge fc1   -> relu1
    edge relu1 -> fc2
    edge fc2   -> relu2
    edge relu2 -> drop
    edge drop  -> fc3
    edge fc3   -> out
}

config {
    batch_size = 64
    device = "cuda"
}
\end{amzcode}
\end{codebox}

Generated PyTorch code:

\begin{codebox}{Generated PyTorch code for MLP}{code-mlp-output}
\begin{amzcode}{python}
import torch
import torch.nn as nn


class MLP(nn.Module):
    def __init__(self):
        super(MLP, self).__init__()
        self.fc1   = nn.Linear(in_features=784, out_features=256)
        self.relu1 = nn.ReLU()
        self.fc2   = nn.Linear(in_features=256, out_features=128)
        self.relu2 = nn.ReLU()
        self.drop  = nn.Dropout(p=0.3)
        self.fc3   = nn.Linear(in_features=128, out_features=10)
        self.out   = nn.Softmax(dim=1)

    def forward(self, x):
        x = self.fc1(x)
        x = self.relu1(x)
        x = self.fc2(x)
        x = self.relu2(x)
        x = self.drop(x)
        x = self.fc3(x)
        x = self.out(x)
        return x


if __name__ == '__main__':
    device = torch.device('cuda')
    model  = MLP().to(device)
    x      = torch.randn(64, 784).to(device)
    print(model(x).shape)
\end{amzcode}
\end{codebox}

% ================================================================
\chapter{Implementation}

\section{Compiler Architecture}

\begin{dfnbox}{Compilation pipeline}{pipeline}
The compiler has six stages:
\begin{enumerate}
    \item \textbf{Parsing}: \code{FileStream} $\to$ \code{NNGraphLexer}
    $\to$ \code{NNGraphParser.program()} $\to$ parse tree
    \item \textbf{Semantic analysis}: \code{ParseTreeWalker} +
    \code{NNGraphCustomListener} $\to$ graph model
    \item \textbf{Topological sort}: compute execution order using
    Kahn's algorithm
    \item \textbf{Shape inference}: propagate tensor shapes through
    the graph and detect dimension mismatches (optional:
    \code{-{}-show-shapes})
    \item \textbf{Code generation}: \code{CodeGenerator.generate()}
    $\to$ Python string
    \item \textbf{Output}: write to \code{.py} file, Graphviz DOT
    output, or ONNX export
\end{enumerate}
\end{dfnbox}

\subsection{File Responsibilities}

\begin{center}
\begin{tabular}{ll}
\toprule
\textbf{File} & \textbf{Role} \\
\midrule
\code{NNGraph.g4}             & Grammar source of truth (56 lines) \\
\code{gen/NNGraphLexer.py}    & Generated by ANTLR4 \\
\code{gen/NNGraphParser.py}   & Generated by ANTLR4 \\
\code{gen/NNGraphListener.py} & Base listener --- should not be edited \\
\code{custom\_listener.py}    & Semantic analysis \\
\code{code\_generator.py}     & PyTorch code generation \\
\code{shape\_inference.py}    & Compile-time tensor shape inference \\
\code{graph\_viz.py}          & Graphviz DOT visualization \\
\code{onnx\_export.py}        & ONNX model export \\
\code{main.py}                & CLI interface and \code{compile\_nng()} \\
\bottomrule
\end{tabular}
\end{center}

\section{Semantic Analyzer}

\begin{dfnbox}{NNGraphCustomListener}{custom-listener}
The class \code{NNGraphCustomListener(NNGraphListener)} in
\code{custom\_listener.py} implements all semantic analysis logic.
Internal state:
\begin{itemize}[noitemsep]
    \item \code{self.nodes}: \code{dict[id, \{type, params, line\}]}
    \item \code{self.edges}: \code{list[\{src, dst, label, line\}]}
    \item \code{self.config}: \code{dict[key, value]}
\end{itemize}
\end{dfnbox}

\subsection{Value Propagation in the Tree}

The \code{exitShape\_expr} and \code{exitValue} methods store values
on the context object itself:

\begin{codebox}{Value propagation via ctx.pvalue in exitShape\_expr}{code-propagation}
\begin{amzcode}{python}
def exitShape_expr(self, ctx):
    ints = [int(ctx.INT_LITERAL(i).getText())
            for i in range(ctx.INT_LITERAL().__len__())]
    ctx.pvalue = tuple(ints)
    ctx.ptype  = tuple
\end{amzcode}
\end{codebox}

This approach ensures that each context, after exiting, has its
children's values available.

\subsection{Validation Rules}

Eight validation rules are executed in two phases:

\textbf{Phase one (during tree walk):}
\begin{enumerate}
    \item \textbf{Duplicate node ID}: in \code{exitNode\_decl}
    \item \textbf{Invalid edge reference}: in \code{exitEdge\_decl}
    \item \textbf{Unknown layer type}: in \code{exitNode\_decl}
    \item \textbf{Parameter type mismatch}: in \code{exitNode\_decl}
\end{enumerate}

\textbf{Phase two (in \code{exitGraph\_block}, after graph completion):}
\begin{enumerate}
    \setcounter{enumi}{4}
    \item \textbf{Undeclared output node}: in
    \code{\_validate\_output\_declared}
    \item \textbf{Orphan node}: in \code{\_validate\_orphans}
    \item \textbf{Residual arity $\neq$ 2}: in
    \code{\_validate\_residual\_arity}
    \item \textbf{Cycle detection via Kahn's algorithm}: in
    \code{\_validate\_no\_cycles}
    \item \textbf{BFS reachability}: in \code{\_validate\_reachability}
\end{enumerate}

\subsection{Error Message Format}

\begin{codebox}{Error format}{code-error-format}
\begin{amzcode}{text}
[line L:C] SemanticError: <message>
[line L] ShapeError at '<node>': <message>
\end{amzcode}
\end{codebox}

\code{L} is the line number (1-based) and \code{C} is the token
column (0-based). All error conditions halt compilation and exit
with code 1.

\subsection{\_FakeCtx for Post-Walk Errors}

Errors raised in \code{exitGraph\_block} do not have a live context.
An internal class \code{\_FakeCtx} carries the line number from the
node's declaration record:

\begin{codebox}{\_FakeCtx class for preserving line numbers}{code-fakectx}
\begin{amzcode}{python}
class _FakeCtx:
    class start:
        line   = info["line"]
        column = 0
\end{amzcode}
\end{codebox}

\subsection{Listener Execution Flow}

The \code{ParseTreeWalker.walk(listener, parse\_tree)} method performs
a depth-first traversal of the parse tree. For each node, it calls
\code{enterX(ctx)} before visiting children and \code{exitX(ctx)}
after all children have been visited.

\begin{dfnbox}{Walker invocation in main.py}{walker-invoke}
\begin{amzcode}{python}
listener = NNGraphCustomListener()
listener.rule_names = parser.ruleNames
walker = ParseTreeWalker()
walker.walk(listener, parse_tree)
\end{amzcode}
\end{dfnbox}

The custom listener overrides only \code{exit*} methods. All
\code{enter*} methods remain as no-ops from the base
\code{NNGraphListener} class.

\subsubsection{Attribute Propagation (Bottom-Up)}

Values flow upward through the tree via context attributes:

\begin{codebox}{exitValue stores parsed value in ctx.pvalue}{code-exitvalue}
\begin{amzcode}{python}
def exitValue(self, ctx):
    if ctx.FLOAT_LITERAL():
        ctx.pvalue = float(ctx.FLOAT_LITERAL().getText())
        ctx.ptype  = float
    elif ctx.INT_LITERAL():
        ctx.pvalue = int(ctx.INT_LITERAL().getText())
        ctx.ptype  = int
    elif ctx.BOOL_LITERAL():
        ctx.pvalue = ctx.BOOL_LITERAL().getText() == 'true'
        ctx.ptype  = bool
    # ... other cases
\end{amzcode}
\end{codebox}

Parent rules read child attributes:

\begin{codebox}{exitParam reads ctx.value().pvalue from child}{code-exitparam}
\begin{amzcode}{python}
def exitParam(self, ctx):
    ctx.param_name  = ctx.ID().getText()
    ctx.param_value = ctx.value().pvalue
    ctx.param_type  = ctx.value().ptype
\end{amzcode}
\end{codebox}

\subsubsection{Semantic Extraction and Detailed Method Behavior}

The custom listener relies on a set of \code{exit*} methods to process the parse tree bottom-up. Below is the detailed explanation of how all listener functions operate:

\begin{itemize}
    \item \textbf{\code{exitShape\_expr(ctx)}}: Parses list of dimensions from \code{INT\_LITERAL} tokens. It builds a Python \code{tuple} and binds it to \code{ctx.pvalue}, setting \code{ctx.ptype} to \code{tuple}.
    \item \textbf{\code{exitValue(ctx)}}: Evaluates literal values (\code{FLOAT\_LITERAL}, \code{INT\_LITERAL}, \code{BOOL\_LITERAL}, \code{STRING}, \code{NONE}, or nested \code{shape\_expr}). It converts the matched text into its corresponding Python representation, attaching the result to \code{ctx.pvalue} and the type to \code{ctx.ptype}.
    \item \textbf{\code{exitModel\_block(ctx)}}: Extracts the neural network model identifier via \code{ctx.ID().getText()} and stores it in \code{self.model\_name}.
    \item \textbf{\code{exitInput\_decl(ctx)}}: Stores the input variable identifier in \code{self.input\_id} and propagates the parsed shape to \code{self.input\_shape}. Additionally, it populates \code{self.nodes} with a pseudo-node of type \code{"\_\_input\_\_"} at the current line number, enabling downstream edge verification to validate connections from the input.
    \item \textbf{\code{exitOutput\_decl(ctx)}}: Saves the output variable identifier in \code{self.output\_id} to track the network's terminal node.
    \item \textbf{\code{exitParam(ctx)}}: Processes key-value parameter configurations. It maps the parameter identifier to \code{ctx.param\_name}, and associates the evaluated values and types to \code{ctx.param\_value} and \code{ctx.param\_type} respectively.
    \item \textbf{\code{exitLayer\_expr(ctx)}}: Processes layer constructor configurations. It identifies the layer class name from the ID and constructs a dictionary mapping parameter names to their parsed metadata. The dictionary is attached to \code{ctx.params}.
    \item \textbf{\code{exitNode\_decl(ctx)}}: Registers nodes and validates structural constraints. It ensures that the node identifier has not been previously declared. It matches the layer type against \code{LAYER\_SCHEMA}, verifies that every declared parameter matches the schema specifications (types and formats), validates that optional quantization annotations (\code{quant}) use a valid mode, and registers the validated node into \code{self.nodes}.
    \item \textbf{\code{exitEdge\_decl(ctx)}}: Processes connections between nodes. It extracts the source and destination identifiers, validates that both endpoints exist in \code{self.nodes}, checks for optional label annotations (e.g., \code{"shortcut"} or \code{"residual\_*"}), and registers the edge in \code{self.edges}.
    \item \textbf{\code{exitConfig\_entry(ctx)}}: Extracts compiler configuration overrides (such as target \code{device} or \code{batch\_size}) and populates the \code{self.config} dictionary.
    \item \textbf{\code{exitGraph\_block(ctx)}}: Triggers after the graph block has been fully walked. It runs five sequential validation procedures: ensuring the output node is declared, verifying that no nodes are orphaned, checking residual input arity, verifying the absence of directed cycles, and validating path reachability from input to output.
    \item \textbf{\code{exitEveryRule(ctx)}}: Intercepts all rule exits. If rule names are loaded, it invokes \code{make\_ast\_subtree} to build and merge abstract syntax tree nodes bottom-up.
    \item \textbf{\code{exitStart(ctx)}}: Finalizes AST construction, binding the complete AST hierarchy under the root rule to the listener's AST collector.
\end{itemize}

\subsubsection{Graph Validation Algorithms}

\code{exitGraph\_block} runs five validation passes after the graph
is fully constructed:

\textbf{1. Orphan Detection} (\code{\_validate\_orphans}):

\begin{codebox}{Orphan detection algorithm}{code-orphan}
\begin{amzcode}{python}
referenced = set()
for e in self.edges:
    referenced.add(e["src"])
    referenced.add(e["dst"])
for node_id, info in self.nodes.items():
    if info["type"] == "__input__":
        continue
    if node_id not in referenced:
        error(f"Node '{node_id}' is declared but never referenced.")
\end{amzcode}
\end{codebox}

Collects all node IDs appearing in edges. Any declared node (except
input) not in this set is orphaned---declared but never connected.

\textbf{2. Cycle Detection} (\code{\_validate\_no\_cycles}):

Uses Kahn's topological sort algorithm. A DAG can be fully sorted;
if any nodes remain with \code{in\_degree > 0} after processing,
a cycle exists.

\begin{codebox}{Kahn's algorithm for cycle detection}{code-kahn}
\begin{amzcode}{python}
in_degree = dict.fromkeys(self.nodes, 0)
adj       = {n: [] for n in self.nodes}
for e in self.edges:
    adj[e["src"]].append(e["dst"])
    in_degree[e["dst"]] += 1

queue   = deque(n for n, d in in_degree.items() if d == 0)
visited = 0
while queue:
    node = queue.popleft()
    visited += 1
    for nxt in adj[node]:
        in_degree[nxt] -= 1
        if in_degree[nxt] == 0:
            queue.append(nxt)

if visited != len(self.nodes):
    remaining = [n for n, d in in_degree.items() if d > 0]
    error(f"Cycle detected involving nodes: {remaining}.")
\end{amzcode}
\end{codebox}

\textbf{3. Reachability Analysis} (\code{\_validate\_reachability}):

Performs BFS from the input node. Every declared node and the output
node must be reachable.

\begin{codebox}{BFS reachability check}{code-bfs}
\begin{amzcode}{python}
adj = {n: [] for n in self.nodes}
for e in self.edges:
    adj[e["src"]].append(e["dst"])

visited = set()
queue   = deque([self.input_id])
while queue:
    node = queue.popleft()
    if node in visited:
        continue
    visited.add(node)
    for nxt in adj[node]:
        queue.append(nxt)

for node_id, info in self.nodes.items():
    if node_id not in visited and info["type"] != "__input__":
        error(f"Node '{node_id}' is not reachable from input.")

if self.output_id not in visited:
    error(f"Output node '{self.output_id}' is not reachable.")
\end{amzcode}
\end{codebox}

Catches disconnected subgraphs and broken input-to-output paths.

\begin{notebox}
Cycle detection and reachability are independent checks. A graph can
be acyclic yet fail reachability (disconnected components), or be
connected yet contain cycles. Both properties are required for a
valid neural network DAG.
\end{notebox}

\section{Code Generator}

\begin{dfnbox}{CodeGenerator}{codegen}
The \code{CodeGenerator} class in \code{code\_generator.py} has
four internal steps:
\begin{enumerate}[noitemsep]
    \item \code{\_compute\_graph()}: build adjacency lists and
    topological sort
    \item \code{\_assign\_vars()}: assign a Python variable name
    to each node
    \item \code{\_emit\_init()} + \code{\_emit\_forward()}: generate
    code lines
    \item \code{\_assemble()}: combine into final Python source
\end{enumerate}
\end{dfnbox}

\section{Variable Assignment Algorithm}

\begin{dfnbox}{Core variable assignment invariant}{var-invariant}
\code{var\_at[n] = 'x'} only when
\code{out\_degree[predecessor] == 1}.

When a node fans out to multiple destinations (\code{out\_degree > 1}),
its tensor is still needed by other branches. Overwriting \code{x}
at this point would corrupt it. Instead, the downstream node receives
its own identifier as a variable name.
\end{dfnbox}

The four complete rules in \code{\_assign\_vars}:

\begin{codebox}{Variable assignment rules in \_assign\_vars}{code-var-rules}
\begin{amzcode}{text}
Rule 1 - Input node:
    var_at[input_id] = 'x'

Rule 2 - Merge node (in_degree > 1):
    var_at[n] = 'x'       if out_degree[n] <= 1
    var_at[n] = node_id   if out_degree[n] > 1

Rule 3 - Branch start (predecessor.out_degree > 1):
    var_at[n] = node_id

Rule 4 - Linear continuation (predecessor.out_degree == 1):
    var_at[n] = var_at[predecessor]
\end{amzcode}
\end{codebox}

\subsection{MLP Trace (Linear Chain)}

\begin{exbox}{Variable trace for MLP}{trace-mlp}
Topology:
\texttt{x -> fc1 -> relu1 -> fc2 -> relu2 -> drop -> fc3 -> out}

Every predecessor has \code{out\_degree == 1} $\Rightarrow$ Rule 4
always applies $\Rightarrow$ all nodes have \code{var\_at = 'x'}.
The \code{forward} method uses \code{x} everywhere.
\end{exbox}

\subsection{ResBlock Trace (Fan-Out Then Merge)}

\begin{exbox}{Variable trace for ResBlock}{trace-resblock}
\code{x} fans out to \code{conv1} and \code{skip} (shortcut).
Therefore \code{out\_degree[x] = 2}.

\begin{itemize}
    \item \code{conv1}: Rule 3 $\Rightarrow$ \code{var\_at = 'conv1'}
    \item \code{bn1, relu1, conv2, bn2}: Rule 4 $\Rightarrow$ all
    \code{'conv1'}
    \item \code{skip} (Residual, \code{in\_degree=2}): Rule 2,
    \code{out\_degree=1} $\Rightarrow$ \code{var\_at = 'x'}
\end{itemize}

Generated code:
\begin{codebox}{Generated forward pass for ResBlock}{code-resblock-forward}
\begin{amzcode}{python}
conv1 = self.conv1(x)
conv1 = self.bn1(conv1)
conv1 = self.relu1(conv1)
conv1 = self.conv2(conv1)
conv1 = self.bn2(conv1)
x     = conv1 + x       # Residual node
x     = self.flatten(x)
\end{amzcode}
\end{codebox}
\end{exbox}

\subsection{InceptionBlock Trace (Two Parallel Branches)}

\begin{exbox}{Variable trace for InceptionBlock}{trace-inception}
\code{out\_degree[x] = 2} (to \code{convA1} and \code{convB1}).

\begin{itemize}
    \item \code{convA1}: Rule 3 $\Rightarrow$ \code{'convA1'}
    \item \code{convB1}: Rule 3 $\Rightarrow$ \code{'convB1'}
    \item \code{convA2}: Rule 4 $\Rightarrow$ inherits \code{'convA1'}
    \item \code{convB2}: Rule 4 $\Rightarrow$ inherits \code{'convB1'}
    \item \code{cat1} (Concat, \code{in\_degree=2}): Rule 2
    $\Rightarrow$ \code{'x'}
\end{itemize}

Generated code:
\begin{codebox}{Generated forward pass for InceptionBlock}{code-inception-forward}
\begin{amzcode}{python}
convA1 = self.convA1(x)
convA1 = self.convA2(convA1)
convB1 = self.convB1(x)
convB1 = self.convB2(convB1)
x      = torch.cat([convA1, convB1], dim=1)  # Concat node
x      = self.bn(x)
x      = self.out(x)
\end{amzcode}
\end{codebox}
\end{exbox}

\subsection{TransformerEncoder Trace (Two Implicit Residuals)}

\begin{exbox}{Variable trace for TransformerEncoder}{trace-transformer}
\code{out\_degree[x] = 2} and \code{out\_degree[norm1] = 2}.

\begin{itemize}
    \item \code{attn}: Rule 3 $\Rightarrow$ \code{'attn'}
    \item \code{drop1}: Rule 4 $\Rightarrow$ \code{'attn'}
    \item \code{norm1} (\code{in\_deg=2, out\_deg=2}): Rule 2
    $\Rightarrow$ \code{'norm1'}
    \item \code{ff1}: Rule 3 $\Rightarrow$ \code{'ff1'}
    \item \code{gelu1, drop2, ff2}: Rule 4 $\Rightarrow$ \code{'ff1'}
    \item \code{norm2} (\code{in\_deg=2, out\_deg=1}): Rule 2
    $\Rightarrow$ \code{'x'}
\end{itemize}

Generated code:
\begin{codebox}{Generated forward pass for TransformerEncoder}{code-transformer-forward}
\begin{amzcode}{python}
attn, _ = self.attn(x, x, x)    # MultiHeadAttn triple-pass
attn     = self.drop1(attn)
norm1    = self.norm1(x + attn)  # implicit residual_1
ff1      = self.ff1(norm1)
ff1      = self.gelu1(ff1)
ff1      = self.drop2(ff1)
ff1      = self.ff2(ff1)
x        = self.norm2(norm1 + ff1) # implicit residual_2
x        = self.out(x)
x        = self.flatten(x)
return x
\end{amzcode}
\end{codebox}
\end{exbox}

\subsection{Special Code Generation Cases}

\begin{center}
\begin{tabular}{llp{7.5cm}}
\toprule
\textbf{Pattern} & \textbf{Condition} & \textbf{Generated Code} \\
\midrule
MultiHeadAttn     & Layer type     & \code{out, \_ = self.attn(x, x, x)} \\
LSTM / GRU        & Layer type     & \code{out, \_ = self.lstm(x)} \\
Residual          & No module      & \code{out = main + shortcut} \\
Add               & No module      & \code{out = a + b} \\
Concat            & No module      & \code{out = torch.cat([a,b], dim=d)} \\
Split             & No module      & \code{split\_0, split\_1, ... = torch.chunk(x,n,dim=d)} \\
Implicit residual & \code{residual\_*} label & \code{out = self.norm(skip + main)} \\
\bottomrule
\end{tabular}
\end{center}

\begin{notebox}
\textbf{Parameter name mapping:} The DSL uses shorter names. The
code generator maps them to PyTorch names:
\code{in\_ch} $\to$ \code{in\_channels},
\code{out\_ch} $\to$ \code{out\_channels},
\code{kernel} $\to$ \code{kernel\_size}.
\end{notebox}

\subsection{Output Code Structure}

\begin{codebox}{Generated output code template}{code-output-template}
\begin{amzcode}{python}
import torch
import torch.nn as nn


class <ModelName>(nn.Module):
    def __init__(self):
        super(<ModelName>, self).__init__()
        # one line per parameterized layer node

    def forward(self, x):
        # one line per node in topological order
        return <output_var>


if __name__ == '__main__':
    device = torch.device('<device>')
    model  = <ModelName>().to(device)
    x      = torch.randn(<batch_size>, <input_shape>).to(device)
    print(model(x).shape)
\end{amzcode}
\end{codebox}

\section{Shape Inference}

\begin{dfnbox}{Shape Inference}{shape-inference}
The \code{shape\_inference.py} module propagates tensor shapes from
the input node along the topological order. For each layer type,
specific rules are applied: Linear changes the last dimension,
Conv2d computes spatial dimensions, Flatten merges dimensions,
and so on.
\end{dfnbox}

When a dimension mismatch is detected, a \code{ShapeError} is
reported with the line number of the node declaration, and
compilation is halted. The \code{-{}-show-shapes} flag displays
each node's output shape.

\section{Graphviz Visualization}

The \code{graph\_viz.py} module produces Graphviz DOT output. Nodes
are color-coded by layer type (green for input, blue for Linear,
orange for Conv, purple for Normalization). The output node is
marked with \code{peripheries=2} and labeled edges are shown with
dashed lines. If shape inference is enabled, shapes are annotated
on the nodes.

\begin{codebox}{Generating and rendering a graph}{code-graphviz}
\begin{amzcode}{bash}
python main.py -i input/inception.nng --graph-viz output/inception.dot
dot -Tpng output/inception.dot -o output/inception.png
\end{amzcode}
\end{codebox}

\section{Training Scaffold}

If the \code{loss} key is set in the \code{config} block, the code
generator produces a complete training loop including loss function,
optimizer, and epoch loop. Supported loss functions:
\code{CrossEntropyLoss}, \code{MSELoss}, \code{BCELoss},
\code{BCEWithLogitsLoss}, \code{L1Loss}, \code{NLLLoss},
\code{SmoothL1Loss}. Supported optimizers:
\code{Adam}, \code{SGD}, \code{AdamW}, \code{RMSprop}.

\section{Quantization Annotations}

The universal \code{quant} parameter can be applied to any node.
Three modes are supported: \code{"dynamic"}, \code{"static"}, and
\code{"qat"}. QAT nodes are wrapped in
\code{torch.quantization.QuantWrapper}. When at least one node has
a quantization annotation, \code{QuantStub} and \code{DeQuantStub}
are added to \code{\_\_init\_\_} and \code{forward}.

\section{ONNX Export}

The \code{onnx\_export.py} module instantiates the generated code
and uses \code{torch.onnx.export} to produce an \code{.onnx} file
with a dynamic batch axis and opset 17.

\begin{codebox}{Exporting a model to ONNX}{code-onnx}
\begin{amzcode}{bash}
python main.py -i input/mlp.nng --onnx output/mlp.onnx
\end{amzcode}
\end{codebox}

% ================================================================
\chapter{Evaluation}


\section{Forward Pass Results}

\begin{exbox}{MLP}{eval-mlp}
\begin{itemize}[noitemsep]
    \item Input: \code{tensor(784)} --- Output: \code{(batch, 10)}
    \item Test: Softmax rows sum to 1 (\code{atol=1e-5})
\end{itemize}
\end{exbox}

\begin{exbox}{ResBlock}{eval-resblock}
\begin{itemize}[noitemsep]
    \item Input: \code{tensor(64, 32, 32)} --- Output:
    \code{(batch, 65536)}
    \item Test: correct shape; deterministic output in \code{eval} mode
\end{itemize}
\end{exbox}

\begin{exbox}{InceptionBlock}{eval-inception}
\begin{itemize}[noitemsep]
    \item Input: \code{tensor(32, 56, 56)} --- Output:
    \code{(batch, 48, 56, 56)}
    \item Branch A: 16 channels; Branch B: 32 channels; merged via
    \code{Concat(dim=1)} $\to$ 48 channels
\end{itemize}
\end{exbox}

\begin{exbox}{TransformerEncoder}{eval-transformer}
\begin{itemize}[noitemsep]
    \item Input: \code{tensor(512, 128)} --- Output:
    \code{(1024, 128)}
    \item Two implicit residual sub-layers without explicit
    \code{Residual()} nodes
    \item \code{MultiHeadAttn} with \code{embed\_dim=128,
    num\_heads=8, batch\_first=True}
\end{itemize}
\end{exbox}

\section{Variable Correctness Proof}

\begin{dfnbox}{Tensor non-overwrite proof}{var-correctness}
For each branching example, the input tensor's \code{data\_ptr()}
was traced through the forward computation:
\begin{itemize}[noitemsep]
    \item \textbf{ResBlock}: \code{x.data\_ptr()} unchanged at the
    point \code{x = conv1 + x}
    \item \textbf{Inception}: \code{x.data\_ptr()} identical when
    \code{convA1(x)} and \code{convB1(x)} are called
    \item \textbf{Transformer}: \code{x.data\_ptr()} unchanged when
    \code{norm1(x + attn)} is called; \code{norm1.data\_ptr()}
    unchanged when \code{norm2(norm1 + ff1)} is called
\end{itemize}

This follows directly from the invariant:
\code{var\_at[n]='x'} only when \code{out\_degree=1}.
\end{dfnbox}

\section{Codebase Statistics}

\begin{center}
\begin{tabular}{lr}
\toprule
\textbf{File/Directory} & \textbf{Lines} \\
\midrule
\code{NNGraph.g4}                    & 56 \\
\code{custom\_listener.py}           & 303 \\
\code{code\_generator.py}            & 327 \\
\code{shape\_inference.py}           & 193 \\
\code{graph\_viz.py}                 & 90 \\
\code{onnx\_export.py}              & 47 \\
\code{main.py}                       & 121 \\
\code{tests/} (4 files)              & 575 \\
\code{required\_code\_collection/}   & $\approx$100 \\
\midrule
\textbf{Total}                        & $\approx$1812 \\
\bottomrule
\end{tabular}
\end{center}

\begin{itemize}[noitemsep]
    \item Language: Python 3.14.3
    \item ANTLR: 4.13.2
    \item PyTorch: 2.10.0
\end{itemize}

% ================================================================
\chapter{Design Decisions}

\section{Listener vs.\ Visitor}

The Listener pattern naturally suits state accumulation. The
\code{exit*} methods can operate directly on \code{self.nodes}
and \code{self.edges} without needing to return values through the
tree.

\section{load\_from\_listener() vs.\ generate(traversal)}

Some compilers use
\code{CodeGenerator.generate(traversal)}, which receives a
post-order token list---suited for stack-based bytecode generation.
For a graph compiler, the natural input is graph structure
(nodes + edges) rather than a linear token stream. Therefore,
\code{load\_from\_listener(listener)} is more direct.

\section{Variable Naming with node\_id}

The language specification examples use single-letter or descriptive
names. This compiler uses the actual node identifier because:
\begin{itemize}[noitemsep]
    \item It is deterministic and works for larger graphs without
    a counter
    \item It is readable---the variable \code{convA1} is clearer
    than \code{a}
    \item The mathematical output is identical
\end{itemize}

% ================================================================
\chapter{Conclusion and Future Work}

\section{Conclusion}

NNGraph DSL provides a six-stage compiler pipeline from \code{.nng}
to executable \code{nn.Module}. The compiler features an ANTLR4
grammar with 14 parser rules, 24 layer types, 8 validation rules,
and a variable assignment algorithm that handles complex branching
without tensor overwrite. The compiler has been validated
against four architectural examples.

\section{Implemented Features}

\begin{center}
\begin{tabular}{lp{10cm}}
\toprule
\textbf{Feature} & \textbf{Description} \\
\midrule
Shape inference        & Tensor shape propagation and mismatch detection at compile time (\code{-{}-show-shapes}) \\
Graphviz visualization & Model graph output in DOT format (\code{-{}-graph-viz}) \\
Training scaffold      & Training loop generation with optimizer and loss from \code{config} block \\
Quantization           & \code{quant} metadata for model quantization conversion \\
ONNX export            & Model export to ONNX format (\code{-{}-onnx}) \\
\bottomrule
\end{tabular}
\end{center}

\section{Future Work}

\begin{center}
\begin{tabular}{lp{10cm}}
\toprule
\textbf{Feature} & \textbf{Description} \\
\midrule
Named sub-graphs   & Reusable macro blocks \\
Dynamic control flow & \code{if}/\code{else} branches within \code{forward} \\
\bottomrule
\end{tabular}
\end{center}

% ================================================================
\chapter*{References}
\markboth{References}{References}
\addcontentsline{toc}{chapter}{References}

\begin{enumerate}
    \item Parr, T. (2013). \textit{The Definitive ANTLR 4 Reference}.
    Pragmatic Bookshelf.
    \item Paszke, A., et al. (2019). PyTorch: An Imperative Style,
    High-Performance Deep Learning Library. \textit{NeurIPS 2019}.
    \item NADER Framework: Neural Architecture Design via
    Representation .
    \item Zhang, A. M. \textit{amznotes LaTeX template}.
    \url{https://github.com/alexmingzhang/latex-notes-template}
\end{enumerate}

\end{document}
